\originalsection{作业 6: 相对速度与质量流}
\begin{exercise}
原题 1. 河水下游速率 1.0 m/s. 学生相对随水漂流的另一船垂直水流以 0.5 m/s 划船, 河宽 800 m. (a) 岸上观察速度大小及与跨河方向的夹角. (b) 下游漂移距离. (c) 横渡时间.
\end{exercise}
\begin{solution}
(a) 速度分量为 $(\text{跨河},\text{下游})=(0.5,1.0)$ m/s, 大小 $\sqrt{1.25}=1.12$ m/s, 源图角 $\tan^{-1}(1.0/0.5)=63.4^\circ$ (与下游方向夹角为 $26.6^\circ$). (c) 横渡时间 $800/0.5=1600$ s. (b) 漂移 1600 m. 河流速度不改变跨河分量.
\end{solution}
\begin{exercise}
原题 2. 零重力中宇航员 1 质量 $m_1$, 初速向右 $v_1$, 宇航员 2 质量 $m_2$, 初速向右 $v_2>v_1$. 1 把质量 $m$ 的球以相对自己速率 $u$ 向运动反方向投出, 2 接住. (a) 求 1 的末速 $v_{f1}$. (b) 求使二人末速相同的 $u$.
\end{exercise}
\begin{solution}
按题面向后投掷, 球速为 $v_{f1}-u$. 投前质量含球, 动量给 $(m_1+m)v_1=m_1v_{f1}+m(v_{f1}-u)$, 所以 (a) $v_{f1}=v_1+mu/(m_1+m)$. 接球后 $v_{f2}=[m_2v_2+m(v_{f1}-u)]/(m_2+m)$. 令两速相等得 (b)
\begin{equation}
u=\frac{m_2(m_1+m)}{m(m_1+m+m_2)}(v_2-v_1).
\end{equation}
图中 2 位于 1 后方并以更大速度追赶, 可接向后投出的球. 捕获是否发生还需初始相对位置合适; 方程仅确定速度条件.
\end{solution}
\begin{exercise}
原题 3. 真空零重力火箭总质量 $2.81\times10^7$ kg, 燃料 $2.46\times10^7$ kg, 干质量 $0.35\times10^7$ kg, 喷气相对速率 3000 m/s. 从静止出发, 总燃烧 510 s. (a) 一级全烧末速. (b) 第一阶段烧 $2.03\times10^7$ kg, 用时 150 s, 求速度增量. (c) 丢弃空罐等 $1.4\times10^6$ kg, 剩干质量 $2.1\times10^6$ kg, 第二段燃料 $4.3\times10^6$ kg, 用时 360 s, 求增量. (d) 求两级末速. (e) 与一级比较并解释.
\end{exercise}
\begin{solution}
每段均用 $\Delta v=u_e\ln(M_i/M_f)$. (a) $3000\ln(2.81/0.35)=6249$ m/s. (b) $3000\ln(2.81/0.78)=3845$ m/s. (c) 丢罐后起始质量 $6.4\times10^6$ kg, $\Delta v_2=3000\ln(6.4/2.1)=3343$ m/s. (d) 合为 7188 m/s. (e) 多约 940 m/s, 因丢掉死重提高后段质量比. 燃烧时间影响推力, 无外力条件下不影响给定质量比的速度增量.
\end{solution}
\begin{exercise}
原题 4. 雨滴从静止下落, 初始质量 $m_0$, 吸收原先静止的云中水分, 增重率 $\dot m=kmv$, $k$ 单位 m$^{-1}$. 忽略阻力. (a) 求加速度方程. (b) 求终端速度.
\end{exercise}
\begin{solution}
向下为正, 输入物料竖直速度为零, 所以 $m\dot v=mg-v\dot m=mg-kmv^2$. (a) $\dot v=g-kv^2$. (b) $v_T=\sqrt{g/k}$. 质量流造成减速项, 不能因为忽略空气阻力就把加速度置为 $g$.
\end{solution}
\begin{exercise}
原题 5. 无动力车干质量 $m_0$, 初速 $v_0$, 光滑水平地面, 竖直落雨以质量流率 $b$ 进入并粘在车上. (a) 在 $[t,t+\Delta t]$ 写速度变化的微分关系. (b) 积分求 $v(t)$.
\end{exercise}
\begin{solution}
$m=m_0+bt$, 雨的水平输入动量为零. (a) $(m_0+bt)\Delta v=-bv(t)\Delta t$ 至一阶. 原题列出的可用变量遗漏 $v(t)$, 这一物理量不可省去. (b) 分离变量给 $v(t)=m_0v_0/(m_0+bt)$.
\end{solution}
